\chapter{Endoscopic Mucosal Resection (EMR)}

\section*{Overview}
Endoscopic mucosal resection (EMR) is an endoscopic technique used to remove superficial lesions of the gastrointestinal tract, such as polyps, adenomas, or early cancers confined to the mucosa and superficial submucosa. It is both a diagnostic and therapeutic procedure because the resected specimen is examined to determine the depth of invasion and completeness of removal.

\section*{Alternative Names}
EMR, endoscopic polypectomy, wide-field endoscopic mucosal resection, submucosal injection polypectomy, cap-assisted mucosal resection

\section*{Principle}
EMR uses a submucosal injection of saline or another solution to lift the lesion away from the underlying muscularis propria, creating a safety cushion. The elevated tissue is then snared and resected with electrocautery, allowing en bloc or piecemeal removal of the lesion. Histologic evaluation of the specimen determines whether the lesion is confined to the mucosa or has invaded the deeper layers. The technique relies on the principle that superficial lesions without deep invasion can be cured with endoscopic resection alone.

\section*{History}
EMR was developed in Japan in the 1980s and 1990s for the treatment of early gastric cancer and later applied to the esophagus, colon, and duodenum. It represents a cornerstone of endoscopic management of early gastrointestinal neoplasia and avoids surgery for many patients.

\section*{Why It Is Done}
EMR is ordered when endoscopy identifies a polyp, adenoma, or early cancer that is suspected to be confined to the superficial layers of the bowel wall. It is used to treat and stage lesions in the esophagus (including Barrett-associated dysplasia), stomach, duodenum, and colorectum. The goal is curative resection of the lesion while avoiding major surgery and its associated morbidity.

\section*{Is This a Primary or Secondary Test or Procedure}
It is a primary therapeutic procedure for lesions deemed endoscopically resectable by careful inspection and optical assessment. It may be secondary when a previously biopsied or incompletely removed lesion requires definitive removal or staging.

\section*{How It Is Performed}
A submucosal injection of normal saline with or without epinephrine and methylene blue is used to lift the lesion and assess for a non-lifting sign. A snare is placed around the lifted tissue, and electrocautery is applied to resect the lesion either en bloc or piecemeal. The specimen is retrieved, oriented, and submitted for histopathologic examination, and resection site clips may be placed to prevent bleeding.

\section*{Preparation}
Colonoscopic EMR requires a full bowel preparation, whereas upper endoscopic EMR requires fasting. Anticoagulant and antiplatelet medications are typically held according to the risk of bleeding, and informed consent is obtained. Imaging and prior histology are reviewed to confirm that the lesion is amenable to endoscopic resection.

\section*{Contraindications and Precautions}
Lesions with deep submucosal invasion, poor lifting, ulceration, or signs of invasive cancer on optical assessment are relative contraindications because of high risk of incomplete resection. Coagulopathy and inability to tolerate sedation or endoscopy preclude the procedure. Large, flat, or circumferential lesions carry higher risks of perforation, bleeding, and recurrence.

\section*{Risks}
Major risks include bleeding, perforation, and postpolypectomy syndrome with pain and fever from transmural thermal injury. Incomplete resection, stricture formation, and local recurrence are additional risks, especially with piecemeal removal of large lesions.

\section*{Aftercare and Recovery}
Patients are monitored for bleeding or perforation and observed for a period appropriate to the site and size of resection. Most resume a normal diet the same day or within a day, with activity restrictions as directed. A repeat endoscopy at 3--6 months is scheduled to assess the resection site for residual or recurrent tissue.

\section*{Outcome and Follow-up}
The pathology report indicates whether the lesion was completely removed, the depth of invasion, and any adverse features such as poor differentiation or lymphovascular invasion. Complete removal of a superficial lesion is considered curative, whereas deep invasion or positive margins requires further endoscopic, surgical, or oncologic management.

\section*{Expected Results}
There is no normal result because EMR treats a lesion; the desired outcome is complete histologic resection of the neoplastic tissue. Success is defined by negative resection margins and absence of deep invasion.

\section*{Follow-up}
Surveillance endoscopy is performed at 3--6 months and then at interval guidelines to detect recurrence at the resection site and metachronous lesions. If histology shows deep invasion or high-risk features, surgical resection with lymphadenectomy may be recommended.

\section*{When to Seek Medical Care}
Follow the procedure team's specific instructions. Seek urgent medical assessment for severe or worsening pain, uncontrolled bleeding, fever or signs of infection, trouble breathing, chest pain, fainting, new weakness or confusion, inability to urinate, or any rapidly worsening symptom. The appropriate warning signs depend on the procedure and the patient's medical condition.

\section*{References}
\begin{enumerate}
  \item MedlinePlus. Endoscopic Mucosal Resection. U.S. National Library of Medicine; 2025.
  \item Current Medical Diagnosis and Treatment. McGraw-Hill; 2025.
\end{enumerate}

\clearpage
